% --- Progression annuelle ---------------------------------------------------
\chapter*{La progression de l'année}
\markboth{La progression de l'année}{La progression de l'année}
\addcontentsline{toc}{chapter}{La progression de l'année}

Le livre est rangé par domaine, parce que c'est ainsi qu'on révise en fin
d'année. Mais on n'enseigne pas dans cet ordre : la Répartition Annuelle du
Programme d'Études répartit les contenus sur cinq périodes, et alterne
volontairement l'analyse, la géométrie et l'algèbre pour ne pas laisser un
domaine dormir six mois.

Le tableau ci-dessous fait la conversion. Lisez-le dans un sens pour préparer
une séance, dans l'autre pour retrouver un chapitre.

\vspace{1em}

\newcommand{\pastille}[2]{\tikz[baseline=-0.2em]{\fill[#1] (0,0) circle (0.24em);}~\small #2}

\begin{center}
\setlength{\tabcolsep}{7pt}
\renewcommand{\arraystretch}{1.45}
\begin{tabular}{@{}p{2.4cm} p{3.3cm} p{8.4cm}@{}}
\toprule
\textsf{\footnotesize\bfseries PÉRIODE} &
\textsf{\footnotesize\bfseries DATES} &
\textsf{\footnotesize\bfseries CHAPITRES TRAVAILLÉS} \\
\midrule

\textsf{\bfseries Première} & \small 8 sept. — 24 oct. &
\pastille{bmtIndigo}{1. Limites et continuité} \newline
\pastille{bmtVetiver}{12. Nombres complexes} \newline
\pastille{bmtIndigo}{2. Dérivabilité} \\

\textsf{\bfseries Deuxième} & \small 3 nov. — 19 déc. &
\pastille{bmtVetiver}{13. Barycentre} \newline
\pastille{bmtVetiver}{14. Transformations du plan \small(première moitié)} \newline
\pastille{bmtIndigo}{3. Étude et représentation graphique} \newline
\pastille{bmtIndigo}{4. Primitives et calcul intégral} \newline
\pastille{bmtLaterite}{11. Suites numériques} \\

\textsf{\bfseries Troisième} & \small 5 janv. — 20 févr. &
\pastille{bmtIndigo}{5. Fonctions logarithmes} \newline
\pastille{bmtIndigo}{6. Exponentielle et fonctions puissances} \newline
\pastille{bmtIndigo}{7. Équations différentielles} \\

\textsf{\bfseries Quatrième} & \small 23 févr. — 27 mars &
\pastille{bmtVetiver}{14. Isométries \small(seconde moitié)} \newline
\pastille{bmtVetiver}{15. Similitudes planes} \newline
\pastille{bmtVetiver}{16. Géométrie dans l'espace} \newline
\pastille{bmtLaterite}{8. Arithmétique dans $\Z$} \\

\textsf{\bfseries Cinquième} & \small 9 avril — 25 juin &
\pastille{bmtLaterite}{9. Systèmes de numération} \newline
\pastille{bmtLaterite}{10. Calcul matriciel} \newline
\pastille{bmtRaphia}{17. Probabilités} \newline
\pastille{bmtRaphia}{18. Variables aléatoires et lois} \\

\bottomrule
\end{tabular}
\end{center}

\vspace{0.6em}

\begin{center}
\small
\pastille{bmtIndigo}{Analyse}\qquad
\pastille{bmtLaterite}{Algèbre}\qquad
\pastille{bmtVetiver}{Géométrie}\qquad
\pastille{bmtRaphia}{Données et probabilités}
\end{center}

\vspace{1em}

\noindent Les dates sont celles de la répartition en vigueur ; elles varient
d'une année à l'autre de quelques jours. Ce qui ne varie pas, c'est l'ordre,
et le fait que la cinquième période porte à la fois les probabilités et les
révisions : prévoyez-la plus courte qu'elle n'en a l'air.
